\documentclass{exam-zh}
\usepackage{edu-practice}

\examsetup{
  question/show-answer = true,
  fillin/show-answer = true,
  solution/show-solution = show-stay,
}

% ---------- 教师版解答样式 ----------
\newtcolorbox{solutionblock}{
  enhanced, breakable,
  colback=edu-blue-bg,
  colframe=edu-blue!25,
  boxrule=0pt,
  borderline west={2pt}{0pt}{edu-blue!50},
  arc=0pt,
  left=3mm, right=3mm, top=2mm, bottom=2mm,
  before skip=2mm,
  after skip=3mm,
  fontupper=\small,
}

% 解答步骤编号
\newcounter{solstep}
\newcommand{\solstep}[2]{%
  \stepcounter{solstep}%
  \par\medskip
  \noindent\hangindent=6mm
  {\color{edu-blue}\bfseries\textcircled{\scriptsize\thesolstep}\enspace #1}\enspace #2\par
}

\begin{document}

\section*{2025学年度第一学期普陀区九年级自主适应练习数学试卷}

\section*{试卷第 1 页}


\needspace{10\baselineskip}
\begin{center}
  \begin{diagramcoltikz}{0.96\linewidth}{}
  \includegraphics[width=\linewidth]{\detokenize{../../../../documents/初三/普陀区-2025学年度第一学期九年级自主适应练习-数学-试卷及参考答案/002.png}}
  \end{diagramcoltikz}
\end{center}


\clearpage
\section*{试卷第 2 页}


\needspace{10\baselineskip}
\begin{center}
  \begin{diagramcoltikz}{0.96\linewidth}{}
  \includegraphics[width=\linewidth]{\detokenize{../../../../documents/初三/普陀区-2025学年度第一学期九年级自主适应练习-数学-试卷及参考答案/003.png}}
  \end{diagramcoltikz}
\end{center}


\clearpage
\section*{试卷第 3 页}


\needspace{10\baselineskip}
\begin{center}
  \begin{diagramcoltikz}{0.96\linewidth}{}
  \includegraphics[width=\linewidth]{\detokenize{../../../../documents/初三/普陀区-2025学年度第一学期九年级自主适应练习-数学-试卷及参考答案/004.png}}
  \end{diagramcoltikz}
\end{center}


\clearpage
\section*{试卷第 4 页}


\needspace{10\baselineskip}
\begin{center}
  \begin{diagramcoltikz}{0.96\linewidth}{}
  \includegraphics[width=\linewidth]{\detokenize{../../../../documents/初三/普陀区-2025学年度第一学期九年级自主适应练习-数学-试卷及参考答案/005.png}}
  \end{diagramcoltikz}
\end{center}


\clearpage
\section*{试卷第 5 页}


\needspace{10\baselineskip}
\begin{center}
  \begin{diagramcoltikz}{0.96\linewidth}{}
  \includegraphics[width=\linewidth]{\detokenize{../../../../documents/初三/普陀区-2025学年度第一学期九年级自主适应练习-数学-试卷及参考答案/006.png}}
  \end{diagramcoltikz}
\end{center}


\clearpage
\section*{试卷第 6 页}


\needspace{10\baselineskip}
\begin{center}
  \begin{diagramcoltikz}{0.96\linewidth}{}
  \includegraphics[width=\linewidth]{\detokenize{../../../../documents/初三/普陀区-2025学年度第一学期九年级自主适应练习-数学-试卷及参考答案/007.png}}
  \end{diagramcoltikz}
\end{center}


\clearpage
\section*{参考答案及解析第 1 页}


\needspace{10\baselineskip}
\begin{center}
  \begin{diagramcoltikz}{0.96\linewidth}{}
  \includegraphics[width=\linewidth]{\detokenize{../../../../documents/初三/普陀区-2025学年度第一学期九年级自主适应练习-数学-试卷及参考答案/008.png}}
  \end{diagramcoltikz}
\end{center}


\clearpage
\section*{参考答案及解析第 2 页}


\needspace{10\baselineskip}
\begin{center}
  \begin{diagramcoltikz}{0.96\linewidth}{}
  \includegraphics[width=\linewidth]{\detokenize{../../../../documents/初三/普陀区-2025学年度第一学期九年级自主适应练习-数学-试卷及参考答案/009.png}}
  \end{diagramcoltikz}
\end{center}


\clearpage
\section*{参考答案及解析第 3 页}


\needspace{10\baselineskip}
\begin{center}
  \begin{diagramcoltikz}{0.96\linewidth}{}
  \includegraphics[width=\linewidth]{\detokenize{../../../../documents/初三/普陀区-2025学年度第一学期九年级自主适应练习-数学-试卷及参考答案/010.png}}
  \end{diagramcoltikz}
\end{center}


\clearpage
\section*{参考答案及解析第 4 页}


\needspace{10\baselineskip}
\begin{center}
  \begin{diagramcoltikz}{0.96\linewidth}{}
  \includegraphics[width=\linewidth]{\detokenize{../../../../documents/初三/普陀区-2025学年度第一学期九年级自主适应练习-数学-试卷及参考答案/011.png}}
  \end{diagramcoltikz}
\end{center}



\end{document}
